\documentclass{article}
\usepackage[T1]{fontenc}
\usepackage{times}
\begin{document}

\section{Introduction}

Paragraph one is set in Times through the T1 encoding, so its glyphs come from a virtual font.
It holds ligatures such as office, affine and difficult, an accented caf\'e and na\"ive, quotes ``like these''
and a dash -- with {\bfseries bold}, {\itshape italic} and {\sffamily sans} words, plus math $\lambda_0 + \sum_j x_j^2$.

Paragraph two is set in Times through the T1 encoding, so its glyphs come from a virtual font.
It holds ligatures such as office, affine and difficult, an accented caf\'e and na\"ive, quotes ``like these''
and a dash -- with {\bfseries bold}, {\itshape italic} and {\sffamily sans} words, plus math $\lambda_1 + \sum_j x_j^2$.

Paragraph three is set in Times through the T1 encoding, so its glyphs come from a virtual font.
It holds ligatures such as office, affine and difficult, an accented caf\'e and na\"ive, quotes ``like these''
and a dash -- with {\bfseries bold}, {\itshape italic} and {\sffamily sans} words, plus math $\lambda_2 + \sum_j x_j^2$.

Paragraph four is set in Times through the T1 encoding, so its glyphs come from a virtual font.
It holds ligatures such as office, affine and difficult, an accented caf\'e and na\"ive, quotes ``like these''
and a dash -- with {\bfseries bold}, {\itshape italic} and {\sffamily sans} words, plus math $\lambda_3 + \sum_j x_j^2$.

Paragraph five is set in Times through the T1 encoding, so its glyphs come from a virtual font.
It holds ligatures such as office, affine and difficult, an accented caf\'e and na\"ive, quotes ``like these''
and a dash -- with {\bfseries bold}, {\itshape italic} and {\sffamily sans} words, plus math $\lambda_4 + \sum_j x_j^2$.

Paragraph six is set in Times through the T1 encoding, so its glyphs come from a virtual font.
It holds ligatures such as office, affine and difficult, an accented caf\'e and na\"ive, quotes ``like these''
and a dash -- with {\bfseries bold}, {\itshape italic} and {\sffamily sans} words, plus math $\lambda_5 + \sum_j x_j^2$.

Paragraph seven is set in Times through the T1 encoding, so its glyphs come from a virtual font.
It holds ligatures such as office, affine and difficult, an accented caf\'e and na\"ive, quotes ``like these''
and a dash -- with {\bfseries bold}, {\itshape italic} and {\sffamily sans} words, plus math $\lambda_6 + \sum_j x_j^2$.

Paragraph eight is set in Times through the T1 encoding, so its glyphs come from a virtual font.
It holds ligatures such as office, affine and difficult, an accented caf\'e and na\"ive, quotes ``like these''
and a dash -- with {\bfseries bold}, {\itshape italic} and {\sffamily sans} words, plus math $\lambda_7 + \sum_j x_j^2$.

Paragraph nine is set in Times through the T1 encoding, so its glyphs come from a virtual font.
It holds ligatures such as office, affine and difficult, an accented caf\'e and na\"ive, quotes ``like these''
and a dash -- with {\bfseries bold}, {\itshape italic} and {\sffamily sans} words, plus math $\lambda_8 + \sum_j x_j^2$.

Paragraph ten is set in Times through the T1 encoding, so its glyphs come from a virtual font.
It holds ligatures such as office, affine and difficult, an accented caf\'e and na\"ive, quotes ``like these''
and a dash -- with {\bfseries bold}, {\itshape italic} and {\sffamily sans} words, plus math $\lambda_9 + \sum_j x_j^2$.

Paragraph eleven is set in Times through the T1 encoding, so its glyphs come from a virtual font.
It holds ligatures such as office, affine and difficult, an accented caf\'e and na\"ive, quotes ``like these''
and a dash -- with {\bfseries bold}, {\itshape italic} and {\sffamily sans} words, plus math $\lambda_10 + \sum_j x_j^2$.

Paragraph twelve is set in Times through the T1 encoding, so its glyphs come from a virtual font.
It holds ligatures such as office, affine and difficult, an accented caf\'e and na\"ive, quotes ``like these''
and a dash -- with {\bfseries bold}, {\itshape italic} and {\sffamily sans} words, plus math $\lambda_11 + \sum_j x_j^2$.

Paragraph thirteen is set in Times through the T1 encoding, so its glyphs come from a virtual font.
It holds ligatures such as office, affine and difficult, an accented caf\'e and na\"ive, quotes ``like these''
and a dash -- with {\bfseries bold}, {\itshape italic} and {\sffamily sans} words, plus math $\lambda_12 + \sum_j x_j^2$.

Paragraph fourteen is set in Times through the T1 encoding, so its glyphs come from a virtual font.
It holds ligatures such as office, affine and difficult, an accented caf\'e and na\"ive, quotes ``like these''
and a dash -- with {\bfseries bold}, {\itshape italic} and {\sffamily sans} words, plus math $\lambda_13 + \sum_j x_j^2$.

\end{document}
